\newpage
\phantomsection
\label{prf:upper-bound-comparison-with-the-supremum}

\begin{remark*}[Return]
\hyperref[prop:upper-bound-comparison-with-the-supremum]{Return to Proposition}
\end{remark*}

\begin{theorem*}[Upper Bound Comparison with the Supremum]
Let $A \subseteq \mathbb{R}$ be non-empty and bounded above, and let $s = \sup A$. For any $u \in \mathbb{R}$, $u$ is an upper bound of $A$ if and only if $s \le u$.
\end{theorem*}

\begin{proof}
\textbf{Professional Standard Proof.}~\\
$(\Rightarrow)$ Suppose $u$ is an upper bound of $A$. Since $s = \sup A$, $s$ is the least upper bound of $A$, so $s \le u$.

$(\Leftarrow)$ Suppose $s \le u$. Since $s$ is an upper bound of $A$, we have $a \le s$ for all $a \in A$. By transitivity of $\le$ on $\mathbb{R}$, $a \le u$ for all $a \in A$. Hence $u$ is an upper bound of $A$.
\end{proof}

\begin{proof}
\textbf{Detailed Learning Proof.}~\\

\textbf{Step 1.} Let $s = \sup A$. Then $s$ is an upper bound of $A$, and for every upper bound $v$ of $A$, one has $s \le v$.

\textbf{Step 2.} $(\Rightarrow)$ Assume $u$ is an upper bound of $A$. Then $u$ belongs to the set of upper bounds of $A$. By the defining property of $s$ as the least upper bound, $s \le u$.

\textbf{Step 3.} $(\Leftarrow)$ Assume $s \le u$. Since $s$ is an upper bound of $A$, $a \le s$ for all $a \in A$.

\textbf{Step 4.} By transitivity of the order on $\mathbb{R}$, from $a \le s$ and $s \le u$, it follows that $a \le u$ for all $a \in A$. Therefore $u$ is an upper bound of $A$.
\end{proof}

\begin{remark*}[Proof structure]
Equivalence proof: each direction follows directly from the definition of supremum and the transitivity of the order on $\mathbb{R}$.
\end{remark*}

\begin{remark*}[Dependencies]~\\
\begin{itemize}
  \item \hyperref[def:supremum]{Definition of supremum}
  \item \hyperref[def:upper-bound]{Definition of upper bound}
  \item Order properties of $\mathbb{R}$
\end{itemize}
\end{remark*}